\section{Rをつかった回帰分析}
\label{b09_Regression_onR}

\subsection{授業内容}
\subsubsection{概要}
回帰分析や重回帰分析について，基本的な概念や用語・推定値の意味は前回までに習得済みであるが，それが統計環境では実際どのようにして算出されるか，算出された数字はどのような意味があるかについて具体的に理解することを目的とする。座学の中では記号としてしか意味を持たなかったものが，実際の数字として得られるとより理解が深まることにもなる。統計環境はRを用いるが，ほかの統計パッケージを用いても同様の出力が得られる。統計環境が変わる可能性もあるので，出力画面の位置で内容を覚えるのではなく，意味的な理解をしておく必要がある。
\subsubsection{コマ主題細目}
\begin{description}

	\item[散布図の描画] 線形回帰はその名の通り線形関係をモデルで表現するものであり，非線形な関係や外れ値の存在などがあれば適切な推定値が得られたことにはならない。分析をする前にまずデータの様子を確認するために，散布図の描画は必ず行わなければならない。Rの基本的なplot関数を用いた散布図の描画方法を学び，データの視覚的な確認方法を理解する。また，複数の変数がある場合の散布図行列(pairs関数)の活用方法についても理解する。
	      \begin{flushright}
		      $\to$ \textcite{kosugi2019}のP.92--98.
	      \end{flushright}

	\item[線形回帰の実行] データが読み込まれれば，$lm$関数によって最小二乗法による回帰分析を行うことができる。Rによる関数関係の記述(formula)の方法を理解し，出力結果から適切な読み取りができるよう，どのような用語でどのような出力がなされているかをしっかり確認する。特に，切片(Intercept)と傾き(係数)の推定値，標準誤差，t値，p値の意味と解釈について理解する。
	      \begin{flushright}
		      $\to$ \textcite{kosugi2018}のP.54--60.
	      \end{flushright}

	\item[回帰診断とQ-Qプロット] lm関数によって得られた回帰モデルの適切性を診断することが重要である。回帰分析の結果オブジェクトをplot関数に渡すことで，残差プロット，Q-Qプロット，スケールロケーションプロットなど，モデル診断に有用なグラフが得られる。特にQ-Qプロットは残差の正規性を確認するための重要なツールであり，理論上の正規分布と実際の残差の分布を比較することで，回帰モデルの前提条件が満たされているかを確認できる。残差プロットでは残差のパターン(漏斗状や系統的な曲線など)を確認し，モデルの妥当性や外れ値の影響を検討する方法を学ぶ。
	      \begin{flushright}
		      $\to$ \textcite{kosugi2018}のP.68--72
	      \end{flushright}

	\item[残差と予測値の確認] $lm$関数から返される結果オブジェクトの中には，予測値$\hat{Y}$がfitted.values，残差がresidualsという変数名で保存されている。これらを取り出したり図示したりすることで，回帰分析の諸特徴を理解する一助となる。数理的な導出は\textcite{kosugi2018}のP.127--135にあるが，具体的な数字としてこれを確認し，理解する。また，residuals.lm()，fitted.lm()などの関数を用いた残差や予測値の取り出し方も学ぶ。
	      \begin{flushright}
		      $\to$ \textcite{kosugi2018}のP.54--60.
	      \end{flushright}

	\item[重回帰分析の実行] 統計環境において重回帰分析を行うのは，説明項を1つ追加するだけで良いので作業としては容易い。ただし解釈に注意が必要であることを再確認し，また標準化偏回帰係数の算出方法についても理解を進める。また変数を増やす前と増やした後で，重相関係数が増加することを確認する。さらに，summary関数から得られる調整済み決定係数(Adjusted R-squared)や，AIC，BICなどのモデル選択指標についても理解する。
	      \begin{flushright}
		      $\to$ \textcite{kosugi2018}のP.63--67.
	      \end{flushright}

\end{description}
\subsubsection{キーワード}
\begin{itemize}
	\item Rの基本plot関数
	\item lm関数とformula表記
	\item Q-Qプロットと残差診断
	\item fitted.valueとresiduals
	\item 標準化偏回帰係数
\end{itemize}
\subsection{授業情報}
\paragraph{コマの展開方法}
演習
\paragraph{標準シラバスにおける位置づけ}
科目番号5 心理学統計法；2統計に関する基礎的な知識；C/D/E エクセル，R，SPSS入門
\subsubsection{予習・復習課題}
\paragraph{予習}
第\ref{b15_useR}講と同じ形式で行われる。RやRStudio，プロジェクトによる管理など基本的なことを改めて教示しないので，前回の復習をかねて空いている時間にPCルームで基本的なR/RStudioの挙動について理解しておくと良い。またRの内部でどのように数字やデータが扱われるかについて，\textcite{kosugi2019}のP.24--36などを参考にみておくと良い。さらに，単回帰分析と重回帰分析の概念的な理解を確認しておくこと。
\paragraph{復習}
授業で扱ったデータを用いて，以下の課題に取り組むこと：

\begin{enumerate}
\item 単回帰分析を実行し，結果を解釈する
\item 回帰診断プロット(特にQ-Qプロット)を描画し，モデルの妥当性を検討する
\item 重回帰分析を実行し，単回帰分析との結果の違いを比較する
\item 残差と予測値を取り出し，散布図上にプロットする
\item 標準化偏回帰係数を算出し，変数間の相対的な影響力を評価する
\end{enumerate}

